\begin{afterword}
\summaryhead

The post opens with Yoda's ``Do, or do not. There is no try.'' The author once took the line for ``Deep Wisdom that is actually quite stupid,'' since success cannot simply be chosen, and now reads it differently.

Saying ``I'm going to flip that switch'' already means making and carrying out a plan that gives the best chance of flipping it. Saying ``I'm going to try to flip that switch'' moves the goal from the switch to having a plan that might flip it. For humans this matters, because it is easy to convince ourselves that we are maximizing our chances, and almost any effort will pass for our hardest if that is all we aim at.

With a quotation from Steven Brust, the post contrasts asking ``What can I do?'', answered from convenient resources, with asking what needs to be done. A lottery ticket bought with one's last dollar is the satirical example. The post ends: ``Don't try your best. Win, or fail. There is no best.''

\responsehead

The post has two parts: an analysis of the word ``try'' and a claim about people. The analysis is clear. If saying one will do a thing already includes trying, then saying one will try to do it can only name a smaller goal: having a plan that might work. The post also concedes the ordinary meaning, that ``try'' marks the possibility of failure, and argues only about what happens when trying becomes the goal.

The force of the argument lies in the claim about people: that a person who aims at having a plan stops checking whether the plan is any good, and that ``Almost any effort will serve'' as one's best. Its core, that ``do your best'' sets a low bar, has outside support. Locke and Latham's review of goal-setting research (2002) reports that specific, difficult goals ``consistently led to higher performance than urging people to do their best.'' Their explanation is close to the post's: do-your-best goals ``have no external referent and thus are defined idiosyncratically,'' so a wide range of performance counts as meeting them. The same review reports that harder goals produced more effort and performance, which fits the post's principle that much of life's challenge is holding oneself to a high enough standard.

The same research also names a limit to the closing advice. ``Win, or fail'' sets the goal on the outcome. On complex tasks that people have not yet learned, an outcome goal can get in the way. In one air-traffic-control simulation, people given an outcome goal learned the task less well and did worse than people told to do their best. In a later study, a specific, difficult learning goal did better than ``do your best.'' On this evidence, the advice fits tasks one already knows how to do better than tasks one must first learn. The post does not discuss this case; it is a limit the research adds, not an error in the post.

In short: a clear argument that aiming to ``try'' lowers the bar, supported in its core by research on goal setting.
\end{afterword}
